\section{Related Work and Contribution Boundary}
\textbf{LLM post-editing.} Raunak et al.~\cite{raunak2023} formalize direct translation post-editing with GPT-4 and report both quality gains and hallucinated edits. Ki and Carpuat~\cite{ki2024} guide LLM post-editing with external error annotations, and xTower~\cite{treviso2024} trains a model to explain and correct translation errors, with expert assessment of its explanations. Our reason-then-edit condition differs: an unmodified editor writes its own brief diagnosis in the same call, and we treat that text as an observable output, not as a verified explanation. LangMark~\cite{velazquez2025} provides human post-edit triplets; FLORES+ references are independent translations, so our reference-based gains are quality proxies and do not measure reduced human editing effort.

\textbf{Exemplar selection.} Vilar et al.~\cite{vilar2023} analyse prompting and example selection for translation. We retrieve one development exemplar with LaBSE~\cite{feng2022}; the similarity is semantic, and one deterministic exemplar is a reproducibility choice, not a claim of optimality.

\textbf{Relationship to the ICAI study.} Ahmadnia et al.~\cite{ahmadnia2026icai} study English-to-German test-time domain adaptation with a frozen OPUS-MT draft, GPT-4 Turbo, multilingual Sentence-BERT top-three exemplars, and an optional second terminology-verification call. Frozen-NMT anchoring and semantic exemplar retrieval are shared with that work and are cited here as established methodology. The present study changes the empirical setting (four Tagalog/Spanish directions, FLORES+), the baseline (sentence-split NLLB-200 3.3B), and the editors (one open-weight and one commercial), uses a single-call reason-then-edit instead of a second verification call, and adds explicit failure and over-editing analysis. Because data, models, and metrics differ, we make no numerical comparison across the two papers.

\textbf{Evaluation.} We use COMET-22~\cite{rei2020,rei2022} as the primary metric, with BLEU~\cite{papineni2002} and chrF++~\cite{popovic2015,popovic2017} reported with full signatures~\cite{post2018}. No scalar metric replaces source inspection, so preservation screens are reported alongside quality.
